% !TEX root = ../../../main.tex
\section{Image}
\label{sec:linear-algebra-i:image}

% Sources: supplied lecture photos and September 24 slides 3 and 12.
\begin{definition}
  The image of \(L:V\to W\) is the set of its possible outputs:
  \[
    \operatorname{im}L=L(V)=\{L(v):v\in V\}\subseteq W.
  \]
  It is a subspace of the codomain \(W\).
\end{definition}

Figure~\ref{fig:linear-algebra-i:kernel-image} shows the kernel and image
inside \(V\) and \(W\). The map sends every vector of \(V\)
to a vector in \(\operatorname{im}L\subseteq W\), and sends every vector
of \(\ker L\subseteq V\) to \(0_W\). The codomain specifies the target
space; the image consists of the outputs actually attained.

\begin{figure}[htbp]
  \centering
  % TikZ adaptation of Pillsmarch's "Kernel and image of linear map.svg".
% Source: https://commons.wikimedia.org/wiki/File:Kernel_and_image_of_linear_map.svg
% Original: Pillsmarch, December 20, 2023. License: CC BY 4.0.
% This redrawing uses grayscale, nested set regions, and two mapping paths.
% Diagram adaptation licensed under https://creativecommons.org/licenses/by/4.0/
\begin{tikzpicture}[
  x=1mm,
  y=1mm,
  font=\footnotesize,
  line width=0.45pt,
  every node/.style={inner sep=2pt},
  map arrow/.style={line width=0.35pt,-{Latex[length=1.6mm,width=1.1mm]}}
]
  % The circles show set containment, not geometric boundaries.
  \draw (19,0) circle[radius=19mm];
  \draw (86,0) circle[radius=20mm];
  \filldraw[fill=black!15] (15,-5) circle[radius=10.5mm];
  \filldraw[fill=black!10] (86,-1) circle[radius=15.5mm];

  \node at (19,26) {Domain \(V\)};
  \node at (86,26) {Codomain \(W\)};
  \node at (52,33) {\(L:V\to W\)};
  \node at (15,-10) {\(\ker L\)};
  \node at (86,-11.5) {\(\operatorname{im}L\)};

  \coordinate (input-vector) at (25,10);
  \coordinate (output-vector) at (83,7);
  \coordinate (kernel-vector) at (15,-2);
  \coordinate (codomain-zero) at (82,-4);

  \draw[map arrow] (input-vector)
    .. controls (43,17) and (66,14) .. (output-vector);
  \draw[map arrow,dashed] (kernel-vector)
    .. controls (43,-6) and (64,-7) .. (codomain-zero);

  \foreach \point in {input-vector,output-vector,kernel-vector,codomain-zero}
    \fill (\point) circle[radius=1.6pt];
  \node at (22,13) {\(v\)};
  \node at (90,6) {\(L(v)\)};
  \node at (15,1) {\(k\)};
  \node at (89,-3) {\(0_W\)};
\end{tikzpicture}

  \caption{Domain, kernel, image, and codomain of \(L:V\to W\).
    The solid line sends \(v\) to \(L(v)\); the dashed line sends
    \(k\in\ker L\) to \(0_W\). Circles show set containment
    schematically, not bounded vector spaces. TikZ redrawing inspired by
    Wikipedia's \emph{Kernel (linear algebra)} page
    \cite{WikipediaKernel} and Pillsmarch's illustration
    \cite{PillsmarchKernelImage},
    licensed under \href{https://creativecommons.org/licenses/by/4.0/}{CC BY 4.0}.}
  \label{fig:linear-algebra-i:kernel-image}
\end{figure}

Recall surjectivity from \S~\ref{sec:appendix:map-properties}.
The codomain need not equal the image. For example, a map into \(\R\)
need not take every real value. The map \(L\) is surjective exactly when
\(\operatorname{im}L=W\). For finite-dimensional \(W\), the subspace
dimension criterion gives
\[
  \dim\operatorname{im}L=\dim W
    \quad\Longleftrightarrow\quad
  \operatorname{im}L=W.
\]
